\documentclass[11pt,letterpaper]{article}

\usepackage[top=.8in, bottom=.8in, left=0.5in, right=0.5in]{geometry}
\usepackage{amsmath,amssymb,amsthm,mathtools}
\usepackage{booktabs}
\usepackage{longtable}
\usepackage{array}
\usepackage{tabularx}
\usepackage{xcolor}
\usepackage{tcolorbox}
\tcbuselibrary{breakable,skins}
\usepackage{enumitem}
\usepackage[T1]{fontenc}
\usepackage{setspace}
\usepackage{fancyhdr}
\usepackage{titlesec}
\usepackage{listings}
\usepackage{courier}
\usepackage{mathrsfs}

\definecolor{RSblue}{RGB}{25,60,120}
\definecolor{RSgreen}{RGB}{20,110,60}
\definecolor{RSamber}{RGB}{180,110,0}
\definecolor{RSgray}{RGB}{90,90,90}

\newtcolorbox{rsbox}{
  enhanced, breakable,
  colback=RSblue!4, colframe=RSblue!60,
  title={\textbf{RS Ontological Statement}},
  fonttitle=\small\bfseries\color{white},
  attach boxed title to top left={yshift=-2mm,xshift=4mm},
  boxed title style={colback=RSblue},
  left=4pt, right=4pt, top=6pt, bottom=4pt}

\newtcolorbox{verifybox}{
  enhanced, breakable,
  colback=RSgreen!4, colframe=RSgreen!50,
  title={\textbf{External Verification (Continuous Lens)}},
  fonttitle=\small\bfseries\color{white},
  attach boxed title to top left={yshift=-2mm,xshift=4mm},
  boxed title style={colback=RSgreen!70!black},
  left=4pt, right=4pt, top=6pt, bottom=4pt}

\newtcolorbox{observebox}{
  enhanced, breakable,
  colback=RSamber!5, colframe=RSamber!60,
  title={\textbf{Structural Observation}},
  fonttitle=\small\bfseries\color{white},
  attach boxed title to top left={yshift=-2mm,xshift=4mm},
  boxed title style={colback=RSamber!80!black},
  left=4pt, right=4pt, top=6pt, bottom=4pt}

\theoremstyle{plain}
\newtheorem{theorem}{Theorem}[section]
\newtheorem{lemma}[theorem]{Lemma}
\newtheorem{proposition}[theorem]{Proposition}
\newtheorem{corollary}[theorem]{Corollary}
\theoremstyle{definition}
\newtheorem{definition}[theorem]{Definition}
\newtheorem{axiom}{Axiom}
\theoremstyle{remark}
\newtheorem{remark}[theorem]{Remark}

%% ---------------------------------------------------------------
%%  BASE ONTOLOGY MACROS
%% ---------------------------------------------------------------
\newcommand{\Z}{\mathbb{Z}}
\newcommand{\SL}{S_{\mathrm{L}}}
\newcommand{\SLe}{S_{\mathrm{L}}^{\mathrm{ext}}}
\newcommand{\SPs}{S_{\mathrm{P}}^{*}}
\newcommand{\IP}{I_{\mathrm{P}}}
\newcommand{\lop}{\lambda}
\newcommand{\aop}{\eta}
\newcommand{\bplus}{\boxplus}
\newcommand{\med}{\oplus}
\renewcommand{\succ}{\operatorname{succ}}
\newcommand{\rank}{\rho}
\newcommand{\Dcross}{\Delta_{\mathrm{cross}}}
\newcommand{\koppa}{\mathbin{\text{\rm\k{o}}}}
\newcommand{\Tcyc}{T_{11}}
\newcommand{\Tbary}{T_{\mathrm{bary}}}
\newcommand{\Tcharge}{\Gamma}
\newcommand{\Dg}{D_{\mathrm{gauge}}}
\newcommand{\Sint}{S_{\mathrm{int}}}
\newcommand{\Sigimb}{\Sigma_{\mathrm{imb}}}
\newcommand{\VS}{\mathcal{V}_S}
\newcommand{\PG}{\mathrm{PG}}
\newcommand{\Gmg}{G_{\mathrm{mg}}}
\newcommand{\Ksat}{K_{\mathrm{sat}}}
\newcommand{\dslip}{\delta_{\mathrm{slip}}}
\newcommand{\Tvac}{T_{\mathrm{vac}}}
\newcommand{\tauX}{\tau_{10}}
\newcommand{\Pithree}{\Pi_3}
\newcommand{\nG}{n_G}
\newcommand{\aSC}{\alpha^{-1}_{\mathrm{int}}}
\newcommand{\piRS}{\pi_{\mathrm{RS}}}
\newcommand{\pa}{\mathfrak{p}_\alpha}
\newcommand{\qa}{\mathfrak{q}_\alpha}
\newcommand{\Da}{\mathcal{D}_\alpha}
\newcommand{\qosc}{\mathfrak{q}_{\mathrm{osc}}}
\newcommand{\RD}{R_D}
\newcommand{\LucN}[1]{L(#1)}
\newcommand{\FibN}[1]{F(#1)}

%% ---------------------------------------------------------------
%%  SHELL DYNAMICS SEMANTIC MACROS
%% ---------------------------------------------------------------
%% Kinematics & Expansion
\newcommand{\KinAmp}{\xi}
\newcommand{\AmpBound}{\upsilon}
\newcommand{\ExcIndex}{\iota}
\newcommand{\SpatDisp}{\varsigma}

%% Dynamics & Inertia
\newcommand{\RestInertia}{\varkappa}
\newcommand{\TotInertia}{\chi}
\newcommand{\KinExcess}{\varepsilon}
\newcommand{\AccAction}{\zeta}
\newcommand{\KinSource}{\Psi}

%% Thermodynamics & Evolution
\newcommand{\ResEns}{\varphi}
\newcommand{\ResEnsAlt}{\vartheta}
\newcommand{\EvolTraj}{\Upsilon}
\newcommand{\RelaxNodes}{\beta}
\newcommand{\GroundEq}{{\omega_{\ast}}}
\newcommand{\CoupledState}{\Phi^{\mathrm{cpl}}}
\newcommand{\StaticState}{\Phi^{\mathrm{st}}}
\newcommand{\GroundProj}{\mathscr{R}}
\newcommand{\InertialSpec}{\mathcal{V}}

\lstset{
  basicstyle=\footnotesize\ttfamily,
  breaklines=true, frame=single,
  language=Python, showstringspaces=false,
  numbers=left, numberstyle=\tiny,
  commentstyle=\color{RSgray},
  keywordstyle=\color{RSblue!70!black}}

\pagestyle{fancy}
\fancyhf{}
\fancyhead[L]{\small\textcolor{RSgray}{RigbySpace Discrete Dynamics}}
\fancyhead[R]{\small\textcolor{RSgray}{Transformative Reciprocal Triadic Structure - Page \thepage}}
\fancyfoot[L]{\small\textcolor{RSgray}{D. Veneziano}}
\fancyfoot[R]{\small\textcolor{RSgray}{Discrete Unification Framework}}
\renewcommand{\headrulewidth}{0.4pt}

\titleformat{\section}
  {\large\bfseries\color{RSblue}}{\thesection}{1em}{}[\titlerule]
\titleformat{\subsection}
  {\normalsize\bfseries\color{RSblue}}{\thesubsection}{1em}{}

\setlength{\parskip}{4pt}
\setlength{\parindent}{0pt}
\onehalfspacing


\begin{document}


\begin{titlepage}
\centering
\vspace*{2cm}
{\Huge\bfseries\color{RSblue}RigbySpace Discrete Dynamics\par}
\vspace{0.5cm}
{\LARGE RigbySpace Shell - Ensemble Transfer Chain\par}
\vspace{0.5cm}
{\large ---\par}
\vspace{2cm}
{\large D.\ Veneziano\par}
\vspace{0.3cm}
{\normalsize RigbySpace Discrete Dynamics\par}
\vspace{0.3cm}
\end{titlepage}

\newpage

\section*{Introduction}

This document records the shell-ensemble transfer chain inside RigbySpace using only RS-native primitives, operators, axioms, and timing structure. The objects used throughout are unreduced ERP states, unreduced elliptic triples, structural tension, vacuum-aligned packets of length \(11\), barycentric shells of length \(33\), and the rest inertia, total inertial load, kinetic excess, and kinematic amplitude ledgers induced on closed shell cycles.

The purpose of the document is to assemble, in a single formal sequence, the exact shell-level transfer law
\[
\TotInertia_M^{\mathrm{cyc}}
=
\RestInertia_M^{\mathrm{cyc}}
+
\AccAction_M^{\mathrm{exc}}(H),
\]
the corresponding coupled/static shell decomposition, the canonical reduction law on admissible support-modulus resonance ensembles, and the ground state equilibrium obtained under rest-inertia-stable admissible ensemble expansion.

No reduced-state quotient, external geometric identification, or non-RS interpretive layer is used. All statements are given as discrete structural definitions, lemmas, theorems, and consequences on the fixed unreduced resonance ensemble.

\section*{I. ERP, Triadic Propagation, and Timing}

\begin{enumerate}
\item \textbf{ERP state}
\[
\mathcal S_L=\{(n,d)\in \mathbb Z^2:d\neq 0\}.
\]
ERPs are never reduced. Structural equality is componentwise equality. Observational equivalence is
\[
(n_1,d_1)\sim(n_2,d_2)\iff n_1d_2=n_2d_1.
\]

\item \textbf{Cross-determinant}
For \(A=(p_a,q_a)\), \(B=(p_b,q_b)\),
\[
\mathcal D_\times(A,B)=p_bq_a-p_aq_b.
\]
Define
\[
W(A,B):=\mathcal D_\times(A,B).
\]
Then
\[
\mathcal D_\times(A,B)=0 \iff A\sim B.
\]

\item \textbf{Width-one bracket}
A discrete convergence bracket is
\[
[L,U]
\]
with
\[
W(L,U)=1.
\]

\item \textbf{Mask condition}
The RS-native object is the width-one integer bracket with oscillation inside it. Continuous irrational appearance is assigned no ontological role.

\item \textbf{Triadic update cycle}
For unreduced pairs \((p,q)\), with phase indexed modulo \(3\),
\[
(p,q)\mapsto(p+q,p),\qquad (p,q)\mapsto(p+q,p),\qquad (p,q)\mapsto(q,p).
\]
This is the exact triadic cycle
\[
\eta,\eta,\psi.
\]

\item \textbf{History is structural}
The third element of propagation is the history element forced by unreducedness.

\item \textbf{Vacuum period}
\[
T_{\mathrm{vac}}=11.
\]

\item \textbf{Triadic interaction period}
\[
N_{\mathrm{tri}}=3.
\]

\item \textbf{Phase slip}
\[
\delta_{\mathrm{slip}}=T_{\mathrm{vac}}\koppa N_{\mathrm{tri}}=11\koppa 3=2.
\]

\item \textbf{Barycentric period}
\[
T_{\mathrm{bary}}=T_{\mathrm{vac}}N_{\mathrm{tri}}=11\cdot 3=33.
\]

\item \textbf{4-4-3 partition}
Over one vacuum cycle,
\[
(N_E,N_M,N_R)=(4,4,3),
\qquad
4+4+3=11.
\]

\item \textbf{Barycentric counts}
Over one barycentric shell,
\[
(N_E,N_M,N_R)_{\mathbb B}=(12,12,9),
\qquad
12+12+9=33.
\]
\end{enumerate}

\section*{II. Unreduced Elliptic State, Packets, and Shells}

\begin{enumerate}[resume]
\item \textbf{Unreduced elliptic state}
\[
S_t=(X_t,Y_t,Z_t)\in\mathbb Z^3,
\]
tracked explicitly without reduction, normalization, or identification of equivalent classes.

\item \textbf{Structural tension}
\[
\tau_t:=X_tZ_{t+1}-X_{t+1}Z_t\in\mathbb Z.
\]

\item \textbf{Vacuum-aligned packet}
\[
\Omega_k=(S_{11k},S_{11k+1},\dots,S_{11k+10}).
\]

\item \textbf{Packet tension}
\[
\Theta_k:=\sum_{j=0}^{10}\tau_{11k+j}.
\]

\item \textbf{Barycentric shell}
\[
\mathbb B_k=(\Omega_{3k},\Omega_{3k+1},\Omega_{3k+2})
=(S_{33k},S_{33k+1},\dots,S_{33k+32}).
\]

\item \textbf{Shell tension}
\[
\Xi_k:=\sum_{j=0}^{32}\tau_{33k+j}.
\]

\item \textbf{Packet decomposition of shell}
\[
\Xi_k=\Theta_{3k}+\Theta_{3k+1}+\Theta_{3k+2}.
\]

\item \textbf{Vacuum alignment}
\[
33\koppa 11=0.
\]

\item \textbf{Triadic alignment}
\[
33\koppa 3=0.
\]

\item \textbf{Minimal simultaneous closure}
\[
11\koppa 3=2\neq 0,\qquad 22\koppa 3=1\neq 0,\qquad 33\koppa 3=0.
\]
Thus \(33\) is the first simultaneous vacuum-triadic closure unit.

\item \textbf{Arithmetic tension ledger}
\[
\mathcal A_0:=0,\qquad \mathcal A_{t+1}:=\mathcal A_t+\tau_t.
\]
Hence
\[
\tau_t=\mathcal A_{t+1}-\mathcal A_t.
\]

\item \textbf{Packet potential}
\[
\mathcal P_k:=\mathcal A_{11k}.
\]

\item \textbf{Packet flux law}
\[
\Theta_k=\mathcal P_{k+1}-\mathcal P_k.
\]

\item \textbf{Shell potential}
\[
\mathcal B_k:=\mathcal A_{33k}=\mathcal P_{3k}.
\]

\item \textbf{Shell law}
\[
\Xi_k=\mathcal B_{k+1}-\mathcal B_k.
\]

\item \textbf{Closed-shell conservation}
If the observed shell sequence closes after \(H\) shells modulo a support modulus \(M\), then
\[
\sum_{k=0}^{H-1}\Xi_k\equiv 0\pmod{M^2}.
\]
\end{enumerate}

\section*{III. Kinematics, Expansion, and Shell Activity}

\begin{enumerate}[resume]
\item \textbf{Kinematic amplitude}
For \(m\in\mathbb Z\),
\[
\KinAmp(0):=0,\qquad
\KinAmp(m):=\min\{r\in\mathbb N:-2^r<m<2^r\}\quad(m\neq 0).
\]

\item \textbf{State kinematic amplitude triple}
\[
\mathfrak r_t:=(\KinAmp(X_t),\KinAmp(Y_t),\KinAmp(Z_t)).
\]

\item \textbf{Packet amplitude boundary}
\[
\AmpBound_X^{\mathrm{pkt}}(k):=\max_{0\le j\le 10}\KinAmp(X_{11k+j}),
\]
\[
\AmpBound_Y^{\mathrm{pkt}}(k):=\max_{0\le j\le 10}\KinAmp(Y_{11k+j}),
\]
\[
\AmpBound_Z^{\mathrm{pkt}}(k):=\max_{0\le j\le 10}\KinAmp(Z_{11k+j}).
\]

\item \textbf{Shell amplitude boundary}
\[
\AmpBound_X^{\mathrm{sh}}(k):=\max_{0\le j\le 32}\KinAmp(X_{33k+j}),
\]
\[
\AmpBound_Y^{\mathrm{sh}}(k):=\max_{0\le j\le 32}\KinAmp(Y_{33k+j}),
\]
\[
\AmpBound_Z^{\mathrm{sh}}(k):=\max_{0\le j\le 32}\KinAmp(Z_{33k+j}).
\]

\item \textbf{Shell amplitude boundary tuple}
\[
\mathfrak G_k^{\mathrm{res}}=
\big(\AmpBound_X^{\mathrm{sh}}(k),\AmpBound_Y^{\mathrm{sh}}(k),\AmpBound_Z^{\mathrm{sh}}(k)\big).
\]

\item \textbf{Cumulative shell amplitude boundary}
\[
\widehat\AmpBound_X(0):=\AmpBound_X^{\mathrm{sh}}(0),\qquad
\widehat\AmpBound_X(k+1):=\max(\widehat\AmpBound_X(k),\AmpBound_X^{\mathrm{sh}}(k+1)),
\]
and similarly for \(\widehat\AmpBound_Y,\widehat\AmpBound_Z\).

\item \textbf{Monotonicity}
For every \(k\),
\[
\widehat\AmpBound_X(k+1)\ge \widehat\AmpBound_X(k),\qquad
\widehat\AmpBound_Y(k+1)\ge \widehat\AmpBound_Y(k),\qquad
\widehat\AmpBound_Z(k+1)\ge \widehat\AmpBound_Z(k).
\]

\item \textbf{Coordinate shell jumps}
\[
\Delta_X^{\mathrm{sh}}(k):=\widehat\AmpBound_X(k+1)-\widehat\AmpBound_X(k),
\]
\[
\Delta_Y^{\mathrm{sh}}(k):=\widehat\AmpBound_Y(k+1)-\widehat\AmpBound_Y(k),
\]
\[
\Delta_Z^{\mathrm{sh}}(k):=\widehat\AmpBound_Z(k+1)-\widehat\AmpBound_Z(k).
\]

\item \textbf{Kinematic-jump triple}
\[
\Delta^{\mathrm{sh}}(k):=
\big(\Delta_X^{\mathrm{sh}}(k),\Delta_Y^{\mathrm{sh}}(k),\Delta_Z^{\mathrm{sh}}(k)\big).
\]

\item \textbf{Kinematic excitation index}
\[
\ExcIndex_{\mathrm{res}}(k):=
\begin{cases}
1,& \Delta_X^{\mathrm{sh}}(k)+\Delta_Y^{\mathrm{sh}}(k)+\Delta_Z^{\mathrm{sh}}(k)>0,\\
0,& \Delta_X^{\mathrm{sh}}(k)+\Delta_Y^{\mathrm{sh}}(k)+\Delta_Z^{\mathrm{sh}}(k)=0.
\end{cases}
\]

\item \textbf{Kinematic excitation criterion}
\[
\ExcIndex_{\mathrm{res}}(k)=1
\iff
\Delta_X^{\mathrm{sh}}(k)>0
\ \text{or}\
\Delta_Y^{\mathrm{sh}}(k)>0
\ \text{or}\
\Delta_Z^{\mathrm{sh}}(k)>0.
\]
\end{enumerate}

\section*{IV. Prime Support, Rest Inertia, Total Inertial Load, and Kinetic Excess}

\begin{enumerate}[resume]
\item \textbf{Support modulus}
Fix a support modulus
\[
M>1.
\]
Let
\[
\Pi(M):=\{p\in\mathbb Z_{>1}\text{ prime}:p\mid M\}.
\]

\item \textbf{Prime shadow state}
For \(p\in\Pi(M)\),
\[
S_t^{(p)}:=(X_t\koppa p,\ Y_t\koppa p,\ Z_t\koppa p).
\]

\item \textbf{Prime divisibility depth of elementary tension}
For \(p\in\Pi(M)\), define
\[
d_p(t):=
\max\{m\in\mathbb Z_{\ge 0}:p^m\mid \tau_t\}
\]
when \(\tau_t\neq 0\), and
\[
d_p(t):=0
\]
when \(\tau_t=0\).

\item \textbf{Elementary prime participation}
\[
c_p(t):=
\begin{cases}
1,& d_p(t)>0,\\
0,& d_p(t)=0.
\end{cases}
\]

\item \textbf{Packet prime load}
\[
D_p^{\mathrm{pkt}}(k):=\sum_{j=0}^{10}d_p(11k+j).
\]

\item \textbf{Packet prime participation}
\[
C_p^{\mathrm{pkt}}(k):=
\begin{cases}
1,& D_p^{\mathrm{pkt}}(k)>0,\\
0,& D_p^{\mathrm{pkt}}(k)=0.
\end{cases}
\]

\item \textbf{Shell prime load}
\[
D_p^{\mathrm{sh}}(k):=\sum_{j=0}^{32}d_p(33k+j).
\]

\item \textbf{Shell prime participation}
\[
C_p^{\mathrm{sh}}(k):=
\begin{cases}
1,& D_p^{\mathrm{sh}}(k)>0,\\
0,& D_p^{\mathrm{sh}}(k)=0.
\end{cases}
\]

\item \textbf{Shell prime-load decomposition}
\[
D_p^{\mathrm{sh}}(k)
=
D_p^{\mathrm{pkt}}(3k)+D_p^{\mathrm{pkt}}(3k+1)+D_p^{\mathrm{pkt}}(3k+2).
\]

\item \textbf{Global shell rest inertia}
\[
\RestInertia_M^{\mathrm{sh}}(k):=\sum_{p\in\Pi(M)}C_p^{\mathrm{sh}}(k).
\]

\item \textbf{Global shell total inertial load}
\[
\TotInertia_M^{\mathrm{sh}}(k):=\sum_{p\in\Pi(M)}D_p^{\mathrm{sh}}(k).
\]

\item \textbf{Shell kinetic excess}
\[
\KinExcess_p^{\mathrm{sh}}(k):=D_p^{\mathrm{sh}}(k)-C_p^{\mathrm{sh}}(k)\in\mathbb Z_{\ge 0}.
\]

\item \textbf{Global shell kinetic excess}
\[
\KinExcess_M^{\mathrm{sh}}(k):=\sum_{p\in\Pi(M)}\KinExcess_p^{\mathrm{sh}}(k).
\]

\item \textbf{Shell transfer identity}
\[
\TotInertia_M^{\mathrm{sh}}(k)
=
\RestInertia_M^{\mathrm{sh}}(k)+\KinExcess_M^{\mathrm{sh}}(k).
\]

\item \textbf{Support exactness}
\[
\RestInertia_M^{\mathrm{sh}}(k)=0
\iff
\TotInertia_M^{\mathrm{sh}}(k)=0
\iff
\KinExcess_M^{\mathrm{sh}}(k)=0.
\]

\item \textbf{Kinetic excess excitation index}
\[
\ExcIndex_{\mathrm{exc}}(k):=
\begin{cases}
1,& \KinExcess_M^{\mathrm{sh}}(k)>0,\\
0,& \KinExcess_M^{\mathrm{sh}}(k)=0.
\end{cases}
\]

\item \textbf{Shell rest inertia excitation index}
\[
\ExcIndex_{\mathrm{rest}}(k):=
\begin{cases}
1,& \RestInertia_M^{\mathrm{sh}}(k)>0,\\
0,& \RestInertia_M^{\mathrm{sh}}(k)=0.
\end{cases}
\]

\item \textbf{Kinetic source spectrum}
\[
\KinSource_M^{\mathrm{exc}}(k)
:=
\{p\in\Pi(M):\KinExcess_p^{\mathrm{sh}}(k)>0\}.
\]

\item \textbf{Spectrum criterion}
\[
\KinSource_M^{\mathrm{exc}}(k)\neq\varnothing
\iff
\KinExcess_M^{\mathrm{sh}}(k)>0.
\]

\item \textbf{Spectrum-rest-inertia lemma}
For every \(p\in\Pi(M)\),
\[
\KinExcess_p^{\mathrm{sh}}(k)>0
\Longrightarrow
C_p^{\mathrm{sh}}(k)=1.
\]

\item \textbf{Excess-to-rest-inertia}
\[
\KinExcess_M^{\mathrm{sh}}(k)>0
\Longrightarrow
\RestInertia_M^{\mathrm{sh}}(k)>0.
\]

\item \textbf{Accumulated action increment}
\[
z_M^{\mathrm{exc}}(k):=\KinExcess_M^{\mathrm{sh}}(k).
\]

\item \textbf{Accumulated action}
\[
\AccAction_M^{\mathrm{exc}}(0):=0,\qquad
\AccAction_M^{\mathrm{exc}}(k+1):=\AccAction_M^{\mathrm{exc}}(k)+z_M^{\mathrm{exc}}(k).
\]

\item \textbf{Accumulated action identity}
For every \(K\ge 1\),
\[
\AccAction_M^{\mathrm{exc}}(K)=\sum_{k=0}^{K-1}\KinExcess_M^{\mathrm{sh}}(k).
\]

\item \textbf{Successive-shell kinetic excess criterion}
\[
\AccAction_M^{\mathrm{exc}}(k+1)>\AccAction_M^{\mathrm{exc}}(k)
\iff
\KinExcess_M^{\mathrm{sh}}(k)>0.
\]
\end{enumerate}

\section*{V. Synthesized RS Admissibility Laws}

\begin{enumerate}[resume]
\item \textbf{Shell-kinetic-excess realization axiom}
For every support modulus \(M\) and every barycentric shell index \(k\),
\[
\KinExcess_M^{\mathrm{sh}}(k)>0
\Longrightarrow
\Delta_X^{\mathrm{sh}}(k)+\Delta_Y^{\mathrm{sh}}(k)+\Delta_Z^{\mathrm{sh}}(k)>0.
\]

\item \textbf{Index form}
\[
\ExcIndex_{\mathrm{exc}}(k)=1
\Longrightarrow
\ExcIndex_{\mathrm{res}}(k)=1.
\]

\item \textbf{Shell-kinematic witness axiom}
For every support modulus \(M\) and every barycentric shell index \(k\),
\[
\Delta_X^{\mathrm{sh}}(k)+\Delta_Y^{\mathrm{sh}}(k)+\Delta_Z^{\mathrm{sh}}(k)>0
\Longrightarrow
\KinExcess_M^{\mathrm{sh}}(k)>0.
\]

\item \textbf{Index form}
\[
\ExcIndex_{\mathrm{res}}(k)=1
\Longrightarrow
\ExcIndex_{\mathrm{exc}}(k)=1.
\]

\item \textbf{Activity equivalence}
For every shell index \(k\),
\[
\ExcIndex_{\mathrm{exc}}(k)=\ExcIndex_{\mathrm{res}}(k).
\]

\item \textbf{Spatial displacement}
\[
\SpatDisp_M^{\mathrm{sh}}(k):=
\Delta_X^{\mathrm{sh}}(k)+\Delta_Y^{\mathrm{sh}}(k)+\Delta_Z^{\mathrm{sh}}(k).
\]

\item \textbf{Exact transfer criterion}
For every shell index \(k\),
\[
\AccAction_M^{\mathrm{exc}}(k+1)>\AccAction_M^{\mathrm{exc}}(k)
\iff
\SpatDisp_M^{\mathrm{sh}}(k)>0.
\]
\end{enumerate}

\section*{VI. Coupled and Static Shell Classes}

\begin{enumerate}[resume]
\item \textbf{Coupled phase state}
\[
\CoupledState_M
:=
\{k:\ExcIndex_{\mathrm{exc}}(k)=1\ \text{and}\ \ExcIndex_{\mathrm{res}}(k)=1\}.
\]

\item \textbf{Static phase state}
\[
\StaticState_M
:=
\{k:\ExcIndex_{\mathrm{exc}}(k)=0\ \text{and}\ \ExcIndex_{\mathrm{res}}(k)=0\}.
\]

\item \textbf{Admissible shell dichotomy}
For every shell index \(k\), exactly one of
\[
k\in\CoupledState_M,
\qquad
k\in\StaticState_M
\]
holds.

\item \textbf{Coupled-shell identification}
\[
\{k:\ExcIndex_{\mathrm{exc}}(k)=1\}
=
\{k:\ExcIndex_{\mathrm{res}}(k)=1\}
=
\CoupledState_M.
\]

\item \textbf{Static-shell identification}
\[
\{k:\ExcIndex_{\mathrm{exc}}(k)=0\}
=
\{k:\ExcIndex_{\mathrm{res}}(k)=0\}
=
\StaticState_M.
\]

\item \textbf{Coupled-shell count}
For \(K\ge 1\),
\[
N_M^{\mathrm{cpl}}(K):=\#(\CoupledState_M\cap\{0,\dots,K-1\}).
\]

\item \textbf{Static-shell count}
For \(K\ge 1\),
\[
N_M^{\mathrm{st}}(K):=\#(\StaticState_M\cap\{0,\dots,K-1\}).
\]

\item \textbf{Count identity}
\[
K=N_M^{\mathrm{cpl}}(K)+N_M^{\mathrm{st}}(K).
\]
\end{enumerate}

\section*{VII. Coupled and Static Sector Inertia}

\begin{enumerate}[resume]
\item \textbf{Coupled-shell rest inertia increment}
\[
z_M^{\mathrm{cpl,rest}}(k)
:=
\begin{cases}
\RestInertia_M^{\mathrm{sh}}(k),& k\in\CoupledState_M,\\
0,& k\notin\CoupledState_M.
\end{cases}
\]

\item \textbf{Coupled-shell kinetic excess increment}
\[
z_M^{\mathrm{cpl,kin}}(k)
:=
\begin{cases}
\KinExcess_M^{\mathrm{sh}}(k),& k\in\CoupledState_M,\\
0,& k\notin\CoupledState_M.
\end{cases}
\]

\item \textbf{Coupled-shell total inertial increment}
\[
z_M^{\mathrm{cpl,tot}}(k)
:=
\begin{cases}
\TotInertia_M^{\mathrm{sh}}(k),& k\in\CoupledState_M,\\
0,& k\notin\CoupledState_M.
\end{cases}
\]

\item \textbf{Static-shell rest inertia increment}
\[
z_M^{\mathrm{st,rest}}(k)
:=
\begin{cases}
\RestInertia_M^{\mathrm{sh}}(k),& k\in\StaticState_M,\\
0,& k\notin\StaticState_M.
\end{cases}
\]

\item \textbf{Static-shell kinetic excess increment}
\[
z_M^{\mathrm{st,kin}}(k)
:=
\begin{cases}
\KinExcess_M^{\mathrm{sh}}(k),& k\in\StaticState_M,\\
0,& k\notin\StaticState_M.
\end{cases}
\]

\item \textbf{Static-shell total inertial increment}
\[
z_M^{\mathrm{st,tot}}(k)
:=
\begin{cases}
\TotInertia_M^{\mathrm{sh}}(k),& k\in\StaticState_M,\\
0,& k\notin\StaticState_M.
\end{cases}
\]

\item \textbf{Coupled rest inertia}
\[
\RestInertia_M^{\mathrm{cpl}}(K):=\sum_{k=0}^{K-1}z_M^{\mathrm{cpl,rest}}(k).
\]

\item \textbf{Coupled kinetic excess}
\[
\KinExcess_M^{\mathrm{cpl}}(K):=\sum_{k=0}^{K-1}z_M^{\mathrm{cpl,kin}}(k).
\]

\item \textbf{Coupled total inertial load}
\[
\TotInertia_M^{\mathrm{cpl}}(K):=\sum_{k=0}^{K-1}z_M^{\mathrm{cpl,tot}}(k).
\]

\item \textbf{Static rest inertia}
\[
\RestInertia_M^{\mathrm{st}}(K):=\sum_{k=0}^{K-1}z_M^{\mathrm{st,rest}}(k).
\]

\item \textbf{Static kinetic excess}
\[
\KinExcess_M^{\mathrm{st}}(K):=\sum_{k=0}^{K-1}z_M^{\mathrm{st,kin}}(k).
\]

\item \textbf{Static total inertial load}
\[
\TotInertia_M^{\mathrm{st}}(K):=\sum_{k=0}^{K-1}z_M^{\mathrm{st,tot}}(k).
\]

\item \textbf{Coupled-sector transfer identity}
\[
\TotInertia_M^{\mathrm{cpl}}(K)
=
\RestInertia_M^{\mathrm{cpl}}(K)+\KinExcess_M^{\mathrm{cpl}}(K).
\]

\item \textbf{Static-sector transfer identity}
\[
\TotInertia_M^{\mathrm{st}}(K)=\RestInertia_M^{\mathrm{st}}(K).
\]

\item \textbf{Static kinetic excess collapse}
\[
\KinExcess_M^{\mathrm{st}}(K)=0.
\]

\item \textbf{Kinetic excess reduction to coupled sector}
\[
\KinExcess_M^{\mathrm{cpl}}(K)=\AccAction_M^{\mathrm{exc}}(K).
\]

\item \textbf{Total rest inertia decomposition}
\[
\sum_{k=0}^{K-1}\RestInertia_M^{\mathrm{sh}}(k)
=
\RestInertia_M^{\mathrm{cpl}}(K)+\RestInertia_M^{\mathrm{st}}(K).
\]

\item \textbf{Total inertial load decomposition}
\[
\sum_{k=0}^{K-1}\TotInertia_M^{\mathrm{sh}}(k)
=
\TotInertia_M^{\mathrm{cpl}}(K)+\TotInertia_M^{\mathrm{st}}(K).
\]

\item \textbf{Total kinetic excess decomposition}
\[
\sum_{k=0}^{K-1}\KinExcess_M^{\mathrm{sh}}(k)
=
\KinExcess_M^{\mathrm{cpl}}(K)=\AccAction_M^{\mathrm{exc}}(K).
\]
\end{enumerate}

\section*{VIII. Global Closed-Cycle Transfer Law}

\begin{enumerate}[resume]
\item \textbf{Closed-cycle rest inertia}
For a closed shell cycle of length \(H\),
\[
\RestInertia_M^{\mathrm{cyc}}:=\sum_{k=0}^{H-1}\RestInertia_M^{\mathrm{sh}}(k).
\]

\item \textbf{Closed-cycle total inertial load}
\[
\TotInertia_M^{\mathrm{cyc}}:=\sum_{k=0}^{H-1}\TotInertia_M^{\mathrm{sh}}(k).
\]

\item \textbf{Closed-cycle transfer law}
\[
\TotInertia_M^{\mathrm{cyc}}
=
\RestInertia_M^{\mathrm{cyc}}+\AccAction_M^{\mathrm{exc}}(H).
\]

\item \textbf{Closed-cycle domination law}
\[
\TotInertia_M^{\mathrm{cyc}}\ge \RestInertia_M^{\mathrm{cyc}}.
\]

\item \textbf{Closed-cycle equality criterion}
\[
\TotInertia_M^{\mathrm{cyc}}=\RestInertia_M^{\mathrm{cyc}}
\iff
\AccAction_M^{\mathrm{exc}}(H)=0.
\]

\item \textbf{Closed-cycle strict-domination criterion}
\[
\TotInertia_M^{\mathrm{cyc}}>\RestInertia_M^{\mathrm{cyc}}
\iff
\AccAction_M^{\mathrm{exc}}(H)>0.
\]

\item \textbf{Shell-count lower bound}
\[
\AccAction_M^{\mathrm{exc}}(H)\ge N_M^{\mathrm{cpl}}(H).
\]

\item \textbf{Strengthened domination law}
\[
\TotInertia_M^{\mathrm{cyc}}
\ge
\RestInertia_M^{\mathrm{cyc}}+N_M^{\mathrm{cpl}}(H).
\]
\end{enumerate}

\section*{IX. Resonance Ensembles, Minima, and Canonical Reduction}

\begin{enumerate}[resume]
\item \textbf{Admissible support-modulus resonance ensemble}
Let
\[
\ResEns
\]
be a nonempty finite set of support moduli such that for every
\[
M\in\ResEns
\]
the observed shell sequence closes after a finite shell length
\[
H_M\in\mathbb Z_{>0}.
\]

\item \textbf{Ensemble rest inertia function}
\[
\RestInertia_{\ResEns}(M):=\RestInertia_M^{\mathrm{cyc}}.
\]

\item \textbf{Ensemble total inertial load function}
\[
\TotInertia_{\ResEns}(M):=\TotInertia_M^{\mathrm{cyc}}.
\]

\item \textbf{Ensemble accumulated action function}
\[
\AccAction_{\ResEns}(M):=\AccAction_M^{\mathrm{exc}}(H_M).
\]

\item \textbf{Ensemble coupled-shell count}
\[
N_{\ResEns}^{\mathrm{cpl}}(M):=N_M^{\mathrm{cpl}}(H_M).
\]

\item \textbf{Moduluswise transfer law}
For every
\[
M\in\ResEns,
\]
\[
\TotInertia_{\ResEns}(M)
=
\RestInertia_{\ResEns}(M)+\AccAction_{\ResEns}(M).
\]

\item \textbf{Moduluswise strengthened domination}
\[
\TotInertia_{\ResEns}(M)\ge \RestInertia_{\ResEns}(M)+N_{\ResEns}^{\mathrm{cpl}}(M).
\]

\item \textbf{Ensemble rest inertia minimum}
\[
\RestInertia_{\ResEns}^{\min}:=\min_{M\in\ResEns}\RestInertia_{\ResEns}(M).
\]

\item \textbf{Ensemble raw total inertial minimum}
\[
\TotInertia_{\ResEns}^{\min}:=\min_{M\in\ResEns}\TotInertia_{\ResEns}(M).
\]

\item \textbf{Rest-inertia-minimizing sub-ensemble}
\[
\ResEns_{\varkappa}
:=
\{M\in\ResEns:\RestInertia_{\ResEns}(M)=\RestInertia_{\ResEns}^{\min}\}.
\]

\item \textbf{Reduced total inertial minimum on rest-inertia-minimizing sub-ensemble}
\[
\TotInertia_{\ResEns}^{\mathrm{red}}
:=
\min_{M\in\ResEns_{\varkappa}}\TotInertia_{\ResEns}(M).
\]

\item \textbf{Canonical support modulus}
Choose
\[
M_{\ResEns}^{\dagger}\in\ResEns_{\varkappa}
\]
such that
\[
\TotInertia_{\ResEns}(M_{\ResEns}^{\dagger})
=
\TotInertia_{\ResEns}^{\mathrm{red}}.
\]

\item \textbf{Canonical accumulated action}
\[
\AccAction_{\ResEns}^{\dagger}
:=
\AccAction_{\ResEns}(M_{\ResEns}^{\dagger}).
\]

\item \textbf{Canonical coupled-shell count}
\[
N_{\ResEns}^{\dagger, \mathrm{cpl}}
:=
N_{\ResEns}^{\mathrm{cpl}}(M_{\ResEns}^{\dagger}).
\]

\item \textbf{Canonical transfer law}
\[
\TotInertia_{\ResEns}^{\mathrm{red}}
=
\RestInertia_{\ResEns}^{\min}
+
\AccAction_{\ResEns}^{\dagger}.
\]

\item \textbf{Canonical strengthened domination}
\[
\TotInertia_{\ResEns}^{\mathrm{red}}
\ge
\RestInertia_{\ResEns}^{\min}
+
N_{\ResEns}^{\dagger, \mathrm{cpl}}.
\]

\item \textbf{Canonical equality criterion}
\[
\TotInertia_{\ResEns}^{\mathrm{red}}
=
\RestInertia_{\ResEns}^{\min}
\iff
\AccAction_{\ResEns}^{\dagger}=0
\iff
N_{\ResEns}^{\dagger, \mathrm{cpl}}=0.
\]

\item \textbf{Canonical strict-domination criterion}
\[
\TotInertia_{\ResEns}^{\mathrm{red}}
>
\RestInertia_{\ResEns}^{\min}
\iff
\AccAction_{\ResEns}^{\dagger}>0
\iff
N_{\ResEns}^{\dagger, \mathrm{cpl}}>0.
\]

\item \textbf{Equilibrium Projection Operator}
\[
\GroundProj(\ResEns):=M_{\ResEns}^{\dagger}.
\]
\end{enumerate}

\section*{X. Admissible Ensemble Expansion}

\begin{enumerate}[resume]
\item \textbf{Admissible ensemble expansion}
For two admissible resonance ensembles
\[
\ResEns,\ResEnsAlt,
\]
define
\[
\ResEns\sqsubseteq_{\mathrm{ens}}\ResEnsAlt
\iff
\ResEns\subseteq\ResEnsAlt.
\]

\item \textbf{Rest inertia minimum under expansion}
If
\[
\ResEns\sqsubseteq_{\mathrm{ens}}\ResEnsAlt,
\]
then
\[
\RestInertia_{\ResEnsAlt}^{\min}\le \RestInertia_{\ResEns}^{\min}.
\]

\item \textbf{Raw total inertial minimum under expansion}
If
\[
\ResEns\sqsubseteq_{\mathrm{ens}}\ResEnsAlt,
\]
then
\[
\TotInertia_{\ResEnsAlt}^{\min}\le \TotInertia_{\ResEns}^{\min}.
\]

\item \textbf{Rest-inertia-stable expansion}
An expansion is rest-inertia-stable if
\[
\RestInertia_{\ResEnsAlt}^{\min}=\RestInertia_{\ResEns}^{\min}.
\]

\item \textbf{Reduced total inertial load under rest-inertia-stable expansion}
If
\[
\ResEns\sqsubseteq_{\mathrm{ens}}\ResEnsAlt
\]
is rest-inertia-stable, then
\[
\TotInertia_{\ResEnsAlt}^{\mathrm{red}}
\le
\TotInertia_{\ResEns}^{\mathrm{red}}.
\]

\item \textbf{Canonical reduction monotonicity}
If
\[
\ResEns\sqsubseteq_{\mathrm{ens}}\ResEnsAlt
\]
is rest-inertia-stable, then
\[
\TotInertia_{\ResEnsAlt}\big(\GroundProj(\ResEnsAlt)\big)
\le
\TotInertia_{\ResEns}\big(\GroundProj(\ResEns)\big).
\]
\end{enumerate}

\section*{XI. Sequential Ground State Stabilization}

\begin{enumerate}[resume]
\item \textbf{Rest-inertia-stable admissible expansion sequence}
A sequence
\[
(\ResEns_n)_{n\in\mathbb Z_{\ge 0}}
\]
is rest-inertia-stable if:
\[
\ResEns_n\sqsubseteq_{\mathrm{ens}}\ResEns_{n+1}
\quad\text{for all }n,
\]
and
\[
\RestInertia_{\ResEns_{n+1}}^{\min}
=
\RestInertia_{\ResEns_n}^{\min}
\quad\text{for all }n.
\]

\item \textbf{Constant rest inertia value}
Define
\[
\RestInertia_\bullet:=\RestInertia_{\ResEns_n}^{\min}
\]
for any \(n\). This is well-defined.

\item \textbf{Reduced total inertial trajectory}
\[
\TotInertia_n:=\TotInertia_{\ResEns_n}^{\mathrm{red}}\in\mathbb Z_{\ge 0}.
\]

\item \textbf{Canonical accumulated action trajectory}
\[
\AccAction_n:=\AccAction_{\ResEns_n}^{\dagger}\in\mathbb Z_{\ge 0}.
\]

\item \textbf{Monotonicity}
For every \(n\),
\[
\TotInertia_{n+1}\le \TotInertia_n.
\]

\item \textbf{Relaxation nodes}
\[
\RelaxNodes:=\{n\in\mathbb Z_{\ge 0}:\TotInertia_{n+1}<\TotInertia_n\}.
\]

\item \textbf{Finite-relaxation lemma}
\[
\RelaxNodes
\]
is finite.

\item \textbf{Ground state index}
\[
N_\ast
:=
\begin{cases}
0,& \RelaxNodes=\varnothing,\\[4pt]
1+\max\RelaxNodes,& \RelaxNodes\neq\varnothing.
\end{cases}
\]

\item \textbf{Ground state stabilization theorem}
For every
\[
n\ge N_\ast,
\]
\[
\TotInertia_n=\TotInertia_{N_\ast}.
\]

\item \textbf{Ground state ensemble}
\[
\ResEns_\ast:=\ResEns_{N_\ast}.
\]

\item \textbf{Ground state total inertial load}
\[
\TotInertia_\ast:=\TotInertia_{\ResEns_\ast}^{\mathrm{red}}.
\]

\item \textbf{Ground state rest inertia}
\[
\RestInertia_\ast:=\RestInertia_{\ResEns_\ast}^{\min}=\RestInertia_\bullet.
\]

\item \textbf{Ground state kinetic excess}
\[
\AccAction_\ast:=\AccAction_{\ResEns_\ast}^{\dagger}.
\]

\item \textbf{Ground state coupled-shell count}
\[
N_\ast^{\mathrm{cpl}}:=N_{\ResEns_\ast}^{\dagger, \mathrm{cpl}}.
\]

\item \textbf{Ground state support modulus}
\[
M_\ast:=\GroundProj(\ResEns_\ast).
\]

\item \textbf{Ground state transfer law}
\[
\TotInertia_\ast=\RestInertia_\ast+\AccAction_\ast.
\]

\item \textbf{Ground state domination law}
\[
\TotInertia_\ast\ge \RestInertia_\ast.
\]

\item \textbf{Ground state strengthened domination}
\[
\TotInertia_\ast\ge \RestInertia_\ast+N_\ast^{\mathrm{cpl}}.
\]

\item \textbf{Ground state equality criterion}
\[
\TotInertia_\ast=\RestInertia_\ast
\iff
\AccAction_\ast=0
\iff
N_\ast^{\mathrm{cpl}}=0.
\]

\item \textbf{Ground state strict-domination criterion}
\[
\TotInertia_\ast>\RestInertia_\ast
\iff
\AccAction_\ast>0
\iff
N_\ast^{\mathrm{cpl}}>0.
\]

\item \textbf{Tail stability}
For every
\[
n\ge N_\ast,
\]
\[
\TotInertia_{\ResEns_n}^{\mathrm{red}}=\TotInertia_\ast,
\qquad
\RestInertia_{\ResEns_n}^{\min}=\RestInertia_\ast,
\qquad
\AccAction_{\ResEns_n}^{\dagger}=\AccAction_\ast.
\]

\item \textbf{Ground state equilibrium}
\[
\GroundEq:=(\RestInertia_\ast,\TotInertia_\ast,\AccAction_\ast).
\]

\item \textbf{Ground state equilibrium law}
\[
\GroundEq=(\RestInertia_\ast,\RestInertia_\ast+\AccAction_\ast,\AccAction_\ast).
\]
\end{enumerate}

\section*{XII. Evolutionary Trajectory Ground State Equilibrium}

\begin{enumerate}[resume]
\item \textbf{Directed admissible rest-inertia-stable evolutionary trajectory}
A directed admissible rest-inertia-stable evolutionary trajectory is a nonempty set
\[
\EvolTraj
\]
of admissible resonance ensembles such that
\[
\forall \ResEns,\ResEnsAlt\in\EvolTraj,\ \exists \phi\in\EvolTraj
\text{ with }
\ResEns\sqsubseteq_{\mathrm{ens}}\phi
\ \text{and}\
\ResEnsAlt\sqsubseteq_{\mathrm{ens}}\phi,
\]
and every expansion inside \(\EvolTraj\) is rest-inertia-stable.

\item \textbf{Trajectory rest inertia value}
For \(\ResEns\in\EvolTraj\), define
\[
\RestInertia_{\EvolTraj}:=\RestInertia_{\ResEns}^{\min}.
\]

\item \textbf{Well-definedness}
\[
\RestInertia_{\EvolTraj}
\]
is well-defined.

\item \textbf{Trajectory inertial spectrum}
\[
\InertialSpec_{\EvolTraj}
:=
\{\TotInertia_{\ResEns}^{\mathrm{red}}:\ResEns\in\EvolTraj\}
\subseteq \mathbb Z_{\ge 0}.
\]

\item \textbf{Trajectory total inertial minimum}
\[
\TotInertia_{\EvolTraj}^{\min}
:=
\min \InertialSpec_{\EvolTraj}.
\]

\item \textbf{Minimizing ensemble}
Choose
\[
\ResEns_{\EvolTraj}^{\min}\in\EvolTraj
\]
such that
\[
\TotInertia_{\ResEns_{\EvolTraj}^{\min}}^{\mathrm{red}}
=
\TotInertia_{\EvolTraj}^{\min}.
\]

\item \textbf{Trajectory accumulated action}
\[
\AccAction_{\EvolTraj}^{\min}
:=
\AccAction_{\ResEns_{\EvolTraj}^{\min}}^{\dagger}.
\]

\item \textbf{Trajectory coupled-shell count}
\[
N_{\EvolTraj}^{\min, \mathrm{cpl}}
:=
N_{\ResEns_{\EvolTraj}^{\min}}^{\dagger, \mathrm{cpl}}.
\]

\item \textbf{Trajectory transfer law}
\[
\TotInertia_{\EvolTraj}^{\min}
=
\RestInertia_{\EvolTraj}+\AccAction_{\EvolTraj}^{\min}.
\]

\item \textbf{Trajectory domination law}
\[
\TotInertia_{\EvolTraj}^{\min}\ge \RestInertia_{\EvolTraj}.
\]

\item \textbf{Trajectory strengthened domination}
\[
\TotInertia_{\EvolTraj}^{\min}\ge \RestInertia_{\EvolTraj}+N_{\EvolTraj}^{\min, \mathrm{cpl}}.
\]

\item \textbf{Trajectory equality criterion}
\[
\TotInertia_{\EvolTraj}^{\min}=\RestInertia_{\EvolTraj}
\iff
\AccAction_{\EvolTraj}^{\min}=0
\iff
N_{\EvolTraj}^{\min, \mathrm{cpl}}=0.
\]

\item \textbf{Trajectory strict-domination criterion}
\[
\TotInertia_{\EvolTraj}^{\min}>\RestInertia_{\EvolTraj}
\iff
\AccAction_{\EvolTraj}^{\min}>0
\iff
N_{\EvolTraj}^{\min, \mathrm{cpl}}>0.
\]

\item \textbf{Cofinal realization theorem}
For every
\[
\ResEnsAlt\in\EvolTraj,
\]
there exists
\[
\phi\in\EvolTraj
\]
such that
\[
\ResEnsAlt\sqsubseteq_{\mathrm{ens}}\phi
\]
and
\[
\RestInertia_{\phi}^{\min}=\RestInertia_{\EvolTraj},
\qquad
\TotInertia_{\phi}^{\mathrm{red}}=\TotInertia_{\EvolTraj}^{\min}.
\]

\item \textbf{Cofinal kinetic excess realization}
For every
\[
\ResEnsAlt\in\EvolTraj,
\]
there exists
\[
\phi\in\EvolTraj
\]
such that
\[
\ResEnsAlt\sqsubseteq_{\mathrm{ens}}\phi
\]
and
\[
\AccAction_{\phi}^{\dagger}
=
\TotInertia_{\EvolTraj}^{\min}-\RestInertia_{\EvolTraj}.
\]

\item \textbf{Trajectory ground state equilibrium}
\[
\GroundEq_{\EvolTraj}
:=
(\RestInertia_{\EvolTraj},\TotInertia_{\EvolTraj}^{\min},\AccAction_{\EvolTraj}^{\min}).
\]

\item \textbf{Trajectory ground state equilibrium law}
\[
\GroundEq_{\EvolTraj}
=
(\RestInertia_{\EvolTraj},\RestInertia_{\EvolTraj}+\AccAction_{\EvolTraj}^{\min},\AccAction_{\EvolTraj}^{\min}).
\]

\item \textbf{Trajectory ground state theorem}
Every directed admissible rest-inertia-stable evolutionary trajectory admits a cofinal ground state equilibrium
\[
\GroundEq_{\EvolTraj}
=
(\RestInertia_{\EvolTraj},\TotInertia_{\EvolTraj}^{\min},\AccAction_{\EvolTraj}^{\min})
\]
such that for every
\[
\ResEnsAlt\in\EvolTraj,
\]
there exists
\[
\phi\in\EvolTraj
\]
with
\[
\ResEnsAlt\sqsubseteq_{\mathrm{ens}}\phi
\]
and
\[
\RestInertia_{\phi}^{\min}=\RestInertia_{\EvolTraj},
\qquad
\TotInertia_{\phi}^{\mathrm{red}}=\TotInertia_{\EvolTraj}^{\min},
\qquad
\AccAction_{\phi}^{\dagger}
=
\TotInertia_{\EvolTraj}^{\min}-\RestInertia_{\EvolTraj}.
\]

\item \textbf{Trajectory closure law}
\[
\TotInertia_{\EvolTraj}^{\min}
=
\RestInertia_{\EvolTraj}
+
\AccAction_{\EvolTraj}^{\min}
\ge
\RestInertia_{\EvolTraj}.
\]

\item \textbf{Ground state equilibrium}
The RS ground state equilibrium for the fixed unreduced resonance ensemble is the cofinal ground state equilibrium
\[
\GroundEq_{\EvolTraj}
=
(\RestInertia_{\EvolTraj},\TotInertia_{\EvolTraj}^{\min},\AccAction_{\EvolTraj}^{\min}).
\]
\end{enumerate}

\end{document}